\section{A logarithmic bridge and an absolute-value no-go}
\label{sec:logarithmic-bridge}

The affine normal form permits one rigorous step through the parity
gate, but only in a weaker averaging regime and without the uniformity
required by the TPC packet.

\subsection{Every fixed fiber cancels logarithmically}

Tao proved the logarithmically averaged two-point Elliott conjecture
\citep{Tao2016}.  In the special case needed here, if
\begin{equation}
 L_1(t)=A_1t+B_1,\qquad L_2(t)=A_2t+B_2,
 \qquad A_1B_2-A_2B_1\ne0,
 \label{eq:fixed-affine-forms}
\end{equation}
are fixed and positive for large \(t\), then
\begin{equation}
 \sum_{t\le T}\frac{\mu(L_1(t))\mu(L_2(t))}{t}
 =o(\log T).
 \label{eq:Tao-mu-correlation}
\end{equation}
This is Corollary~5 of \citet{Tao2016}, applied with
\(g_1=g_2=\mu\) and \(\omega(T)=T\).  The M\"obius function satisfies
the corrected strong nonpretentiousness condition used there; Tao
deduces this from the classical zero-free region.  A fixed initial
shift may be made first if one of the constants \(B_i\) is negative,
so both forms lie in the positive integers throughout the sum.

\begin{proposition}[Log-weighted cancellation on a fixed fiber]
\label{prop:fixed-fiber-log}
Under \eqref{eq:fixed-affine-forms},
\begin{equation}
 \boxed{
 \sum_{t\le T}\frac{a(L_1(t))a(L_2(t))}{t}
 =o((\log T)^3).}
 \label{eq:fixed-fiber-log}
\end{equation}
In particular, every fixed fiber in
\cref{thm:affine-complement-fiber} has logarithmically averaged
cancellation.
\end{proposition}

\begin{proof}
Set
\[
 A(x)=\sum_{t\le x}\frac{\mu(L_1(t))\mu(L_2(t))}{t}
 =o(\log x)
\]
by \eqref{eq:Tao-mu-correlation}.  The weight
\(w(t)=\log L_1(t)\log L_2(t)\) satisfies
\[
 w(t)\ll(\log(2t))^2,
 \qquad w'(t)\ll\frac{\log(2t)}t
\]
away from a fixed initial interval.  Partial summation gives
\[
 A(T)w(T)-\int_2^T A(t)w'(t)\,dt=o((\log T)^3).
\]
The two minus signs in \(a(L_1)a(L_2)\) cancel.
\end{proof}

For a fixed complement cutoff \(Z\), define
\begin{equation}
 \cT_{m}^{[Z]}(j)
 =\sum_{\substack{r\mid mj+h\\r\le Z}}a((mj+h)/r).
 \label{eq:fixed-complement-truncation}
\end{equation}

\begin{corollary}[Finite complement reassembly]
\label{cor:finite-complement-log}
Fix \(h\ne0\), distinct positive rows \(m_1,m_2\) coprime to \(h\),
and \(Z\ge1\).  Then
\begin{equation}
 \boxed{
 \sum_{j\le J}\frac{
  \cT_{m_1}^{[Z]}(j)\cT_{m_2}^{[Z]}(j)}j
 =o_{m_1,m_2,h,Z}((\log J)^3).}
 \label{eq:finite-complement-log}
\end{equation}
\end{corollary}

\begin{proof}
Expand the two finite divisor sums.  For each fixed \(r,s\le Z\), the
two divisibility conditions are either incompatible or select one
class \(j=j_0+[r,s]t\).  In the latter case,
\cref{thm:affine-complement-fiber} gives two affine quotient forms with
nonzero determinant.  Also
\[
 \frac1{j_0+[r,s]t}
 =\frac1{[r,s]t}+O_{r,s}(t^{-2}).
\]
Apply \cref{prop:fixed-fiber-log}; the \(t^{-2}\) error converges after
the logarithmic weights.  There are only finitely many pairs.
\end{proof}

The fixedness can be relaxed only ineffectively by a diagonal argument.

\begin{corollary}[An arbitrarily slow moving window]
\label{cor:slow-window}
For fixed \(h\ne0\), there exists a nondecreasing function
\(K(J)\to\infty\) such that
\begin{equation}
 \max_{\substack{1\le m_1,m_2,Z\le K(J)\\m_1\ne m_2\\
                  (m_1m_2,h)=1}}
 \frac1{(\log J)^3}
 \left|\sum_{j\le J}\frac{
 \cT_{m_1}^{[Z]}(j)\cT_{m_2}^{[Z]}(j)}j\right|
 \longrightarrow0.
 \label{eq:slow-window}
\end{equation}
\end{corollary}

\begin{proof}
For each integer \(K\), \cref{cor:finite-complement-log} holds
uniformly over the finite set of triples \(m_1,m_2,Z\le K\), with the
maximum over an empty set interpreted as zero.  Choose
strictly increasing thresholds \(J_K\ge K\) beyond which the normalized
maximum is at most \(1/K\), and set
\(K(J)=\max\{K:J_K\le J\}\), with \(K(J)=1\) below \(J_1\).
\end{proof}

This is a genuine moving-parameter consequence, but it supplies no
effective rate at all.  It is many orders of magnitude weaker than
the polynomial scales \(m_i\asymp Q=X^\beta\) and
\(r,s\ll X/S\) in \eqref{eq:exact-two-fiber-reassembly}.

\subsection{Why absolute values cannot replace the correlation}

We next show that the unit-complement fiber is large after signs are
removed.  This is a componentwise obstruction, not a lower bound for
the complete signed residual packet.

For fixed primitive rows, let \(\nu_p\) be the number of solutions
\(j\pmod{p^2}\) of
\begin{equation}
 (m_1j+h)(m_2j+h)\equiv0\pmod{p^2}.
 \label{eq:nu-p}
\end{equation}
Each form has at most one root modulo \(p^2\), and it has none if
\(p\mid m_i\), because then \(p\nmid h\).  Hence \(0\le\nu_p\le2\),
and
\begin{equation}
 \mathfrak S_{\rm sf}(m_1,m_2;h)
 =\prod_p\left(1-\frac{\nu_p}{p^2}\right)>0.
 \label{eq:squarefree-singular-series}
\end{equation}

\begin{lemma}[Simultaneous squarefree density]
\label{lem:simultaneous-squarefree}
For fixed \(m_1\ne m_2\) with \((m_1m_2,h)=1\),
\begin{equation}
 \sum_{j\le J}
 \mu^2(m_1j+h)\mu^2(m_2j+h)
 =\mathfrak S_{\rm sf}(m_1,m_2;h)J+o(J).
 \label{eq:simultaneous-squarefree}
\end{equation}
\end{lemma}

\begin{proof}
For a fixed prime bound \(P\), the Chinese remainder theorem gives the
exact density
\(\prod_{p\le P}(1-\nu_p/p^2)\) of classes avoiding all square
divisors \(p^2\) with \(p\le P\), up to \(O_P(1/J)\).  The number of
\(j\le J\) for which one of the two forms is divisible by \(p^2\) for
some \(p>P\) is
\[
 \ll\sum_{P<p\ll\sqrt J}\left(\frac J{p^2}+1\right)
 =O(J/P)+O(\sqrt J/\log J),
\]
where the implied constants may depend on the fixed rows and \(h\).
Let \(J\to\infty\) and then \(P\to\infty\).  Absolute convergence and
positivity of
\eqref{eq:squarefree-singular-series} follow from \(\nu_p\le2\).
\end{proof}

\begin{proposition}[Absolute-mass no-go]
\label{prop:absolute-nogo}
For the unit-complement coefficient in
\eqref{eq:unit-parity-coefficient},
\begin{equation}
 \boxed{
 \sum_{J/2<j\le J}
 \left|\mu(m_1j+h)\mu(m_2j+h)\right|
 \log(m_1j+h)\log(m_2j+h)
 \gg_{m_1,m_2,h}J(\log J)^2.}
 \label{eq:absolute-nogo}
\end{equation}
Therefore a fiberwise triangle inequality or sign-blind estimate cannot
make this fixed, unweighted component power-saving.
\end{proposition}

\begin{proof}
By \cref{lem:simultaneous-squarefree}, a positive proportion of the
integers in \((J/2,J]\) make both forms squarefree.  On that interval
both logarithms are \(\asymp\log J\), and
\(|\mu(n)|=\mu^2(n)\).
\end{proof}

There is no contradiction between
\cref{prop:fixed-fiber-log,prop:absolute-nogo}: the first retains the
product of the two M\"obius signs and uses logarithmic averaging; the
second deliberately removes exactly that cancellation.
